\section{Three Levers That Are Not Levers}
\label{sec:levers}

When a model underperforms, the reflex is to make it bigger, give it a finer grid, or feed it more
inputs. We tried all three. None of them is the bottleneck here, and we report that because a
negative result about the obvious fix saves other people the same weeks.

\subsection{Capacity is not the bottleneck}
On held-out floods with terrain only, CSI is 0.221 at 2.0\,M parameters, 0.232 at 7.9\,M, 0.221 at
31.4\,M and 0.219 at 70.6\,M. That is flat across a 35-fold range and entirely inside seed noise.

\subsection{Resolution is not the bottleneck either}
Our standing explanation for poor co-location was that real ponds are around 810\,m$^2$, roughly
28.5\,m across, and were being scored on 60\,m cells. Section~\ref{sec:twin} is where that 810\,m$^2$
figure comes from, and it is a good explanation for the \emph{magnitude} error. It is not the
explanation for the location error. Re-running the whole experiment at 30\,m quadruples the cell
count and changes nothing. The skill multiple over all-wet is 10.4 at 30\,m against 9.7 at 60\,m.
Raw CSI is not comparable across grids because the base rates differ, which is why the multiple is
the column to read.

\subsection{Rainfall is not usable, and it actively hurts}
This is the most counter-intuitive result in the paper, so we tested it four ways.

\textbf{Ablation.} Removing the rainfall channels \emph{improves} the model's ranking quality by
0.124, at about seven seed standard deviations, with the three-seed ranges disjoint.

\textbf{Permutation control.} We shuffled which rainfall vector goes with which storm, preserving
every rainfall distribution exactly and destroying only the pairing. The result lands monotonically
between real rain and no rain, at 0.184 then 0.272 then 0.313. Real rainfall is strictly the worst
of the three inputs. The network fits a rain-to-wetness mapping on the training storms that does
not transfer.

\textbf{Two alternative explanations, tested and refuted.} The first was tidal contradiction
between the two Mumbai tiles, which predicted that dropping the tidal tile would stop the harm. It
did not. The second was failure to extrapolate to larger storms, but rain costs 0.09 CSI within the
training distribution too.

\textbf{The boundary.} On an unseen region, rainfall neither helps nor hurts, at 0.1460 with rain
against 0.1474 without, replicated at 15 areas. So the claim is not that rainfall is universally
harmful. It is that rainfall lets the model take a shortcut on a city it has already seen, and that
shortcut does not survive the move to new ground.

\subsection{Why there is nothing in the rainfall to extract}
Two properties of the free rainfall archive explain the result. Nine probe points spanning a 15\,km
tile collapse to a single archive cell, so spatial rainfall is not merely unused. It is
unavailable.

And the source was not what we previously stated. Open-Meteo's archive endpoint, queried without an
explicit model parameter, returns ECMWF IFS at 9\,km rather than ERA5. Over the same window and
point that is 554.2\,mm against ERA5's 1\,213.4\,mm, a factor of 2.19 which exceeds any effect our
physics calibration ever produced. The twin was never calibrated on ERA5. We now pin the model
explicitly, which reproduces every number in this paper exactly. This is the third correction we
make against our own earlier statements.

\section{Why the Physics Router Fails}
\label{sec:dem}

Sections \ref{sec:twin} and \ref{sec:levers} say that terrain contains the information and that the
solver cannot use it. This section says why, and the answer is a property of the elevation data
rather than of the solver.

We re-simulate under a ten-member ensemble of elevation perturbations that are spatially
correlated, with a standard deviation of 1\,m and a correlation length of five cells, which is
FABDEM's nominal error class \cite{marsh2024}. Spatially correlated means neighbouring cells are
wrong together rather than independently, which is how real elevation error behaves.

Total flooded area is robust. Patna gives $10.79 \pm 0.35$\,km$^2$ at 100\,mm. But only
0.58\,km$^2$, which is 5\,\%, floods in at least 90\,\% of the ensemble members. In Bengaluru, with
ten times steeper relief, 4.56 of 8.82\,km$^2$, or 52\,\%, is robust. Fig.~\ref{fig:unc} shows it.

\begin{figure}[!t]
\centering
\includegraphics[width=0.86\linewidth]{figures/flood_uncertainty.png}
\caption{Elevation ensemble at $\pm1$\,m, showing flood probability over ten members for Patna at
100\,mm. The total area is stable but its location moves. Only 5\,\% of flooded area is common to
at least 90\,\% of members on these flat deltaic plains, against 52\,\% in hillier Bengaluru.}
\label{fig:unc}
\end{figure}

Calibrating 16 per-land-class roughness and infiltration multipliers by gradient descent moves
held-out CSI only from 0.0483 to 0.0488, with the multipliers staying between 0.94 and 1.14.
Roughness and infiltration change \emph{how deep} the water is, not \emph{where} it is. We also
ruled out a georeferencing bug, because mirroring the radar mask doubles CSI, which is alarming
until you check that the elevation model is provably correctly oriented. The correlation between
negative elevation and JRC water is 0.80 unflipped against 0.05 flipped, and the Ganga sits 7.8\,m
below the domain mean. The flip is a coincidental mirror and is not applied.

The scope of this diagnosis is narrower than we previously claimed. On a flood plain where
metre-scale elevation noise exceeds the relief that decides which street ponds, a solver that
routes water downhill cannot place it reliably, and its calibration gradient with respect to
location is essentially zero. That is why the twin sits at 0.054. It is \emph{not} a proof that the
terrain data lacks the information, because a model reading relative micro-relief from the same
rasters reaches 0.288. The correct statement is that absolute elevation is too noisy to route over,
while relative elevation is not too noisy to learn from.

Two things follow. Reported skill in this regime should always come with both such an ensemble and
the trivial baselines of Table~\ref{tab:trivial}. And planning tools should be judged on decisions
that survive this uncertainty, which is what we turn to next.
